\section{Косвенная адресация}

\subsection{Что такое косвенная адресация и для чего она применяется?}
\begin{definition}[Косвенная адресация]
\textit{Indirect addressing} — это способ доступа к объекту не по его прямому имени, а через некоторую промежуточную сущность: указатель, ссылку, итератор, \texttt{string\_view} и т.п.
\end{definition}

Косвенная адресация нужна, когда:
\begin{itemize}
    \item объект должен быть доступен не только по своему имени;
    \item нужно передать объект в функцию без копирования;
    \item нужно работать с «чужим» объектом через абстракцию;
    \item нужно обойти коллекцию или часть коллекции;
    \item требуется отделить интерфейс от конкретного хранения данных.
\end{itemize}

\subsection{Какие механизмы языка C++ и компоненты стандартной библиотеки вовлекаются при косвенной адресации?}
Основные механизмы:
\begin{itemize}
    \item указатели;
    \item ссылки;
    \item итераторы;
    \item \texttt{std::string\_view};
    \item в более широком смысле — умные указатели и обёртки над ссылками.
\end{itemize}

\begin{remark}
Все эти средства выполняют одну идею: дать доступ к объекту через другое имя или через позицию, не копируя сам объект без необходимости.
\end{remark}

\subsection{Какое количество объектов как правило задействовано при косвенной адресации: при использовании указателя? ссылки? итератора?}
Обычно:
\begin{itemize}
    \item \textbf{указатель} связывает два объекта: сам указатель и объект, на который он указывает;
    \item \textbf{ссылка} также связывает два объекта: имя-ссылку и объект-цель;
    \item \textbf{итератор} обычно тоже связан с двумя сущностями: самим итератором и элементом коллекции, на который он указывает.
\end{itemize}

\subsection{Косвенная адресация и константность. Один объект — две разновидности константности. Два объекта — четыре разновидности константности}
Когда есть один объект и одна косвенная сущность, можно рассматривать две независимые характеристики:
\begin{itemize}
    \item можно ли изменять \textbf{объект} через это имя;
    \item можно ли изменять саму \textbf{косвенную сущность} как объект языка.
\end{itemize}

Если рассматривать пару «объект + косвенность», то возникают четыре комбинации:
\begin{enumerate}
    \item объект изменяем, косвенность изменяема;
    \item объект неизменяем, косвенность изменяема;
    \item объект изменяем, косвенность неизменяема;
    \item объект неизменяем, косвенность неизменяема.
\end{enumerate}

\begin{remark}
Для указателей и итераторов эти варианты хорошо видны. Для ссылок картина беднее, потому что ссылка в C++ не перепривязывается.
\end{remark}

\subsection{Указатели: «левоконстантный», «правоконстантный» и «дваждыконстантный» указатель}
Пусть есть указатель \texttt{T* p}. Тогда константность может относиться:
\begin{itemize}
    \item к тому, можно ли менять \textbf{содержимое}, на которое указывает \texttt{p};
    \item к тому, можно ли менять \textbf{сам указатель} \texttt{p}.
\end{itemize}

\begin{definition}[Левоконстантный указатель]
\texttt{const T* p} или \texttt{T const* p} — указатель на константный объект.  
Нельзя менять объект через \texttt{p}, но можно переназначать сам указатель.
\end{definition}

\begin{definition}[Правоконстантный указатель]
\texttt{T* const p} — сам указатель константен.  
Можно менять объект через \texttt{p}, но нельзя переназначить \texttt{p} на другой объект.
\end{definition}

\begin{definition}[Дваждыконстантный указатель]
\texttt{const T* const p} — и сам указатель, и объект под ним константны.  
Нельзя ни менять объект через \texttt{p}, ни переназначать \texttt{p}.
\end{definition}

\subsection{Ссылки: константность и неконстантность}
\begin{definition}[Неконстантная ссылка]
\texttt{T\& r} — ссылка, через которую можно менять объект.
\end{definition}

\begin{definition}[Константная ссылка]
\texttt{const T\& r} — ссылка только для чтения: через неё нельзя изменять объект.
\end{definition}

Для ссылок важное отличие от указателей такое:
\begin{itemize}
    \item ссылка не может быть пустой;
    \item обычную ссылку нельзя перепривязать;
    \item поэтому у ссылки нет полноценной аналогии с «правоконстантным» указателем.
\end{itemize}

\begin{remark}
В разговорной модели можно говорить, что у ссылки есть константность относительно объекта, на который она смотрит, но сама ссылка как «переменная-сущность» не ведёт себя как указатель.
\end{remark}

\subsection{Итераторы: \texttt{const\_iterator}, \texttt{const iterator} и \texttt{const const\_iterator}}
У контейнеров STL обычно есть:
\begin{itemize}
    \item \texttt{iterator} — итератор, через который можно менять элемент;
    \item \texttt{const\_iterator} — итератор, через который можно только читать элемент;
    \item \texttt{const iterator} — константный объект типа \texttt{iterator};
    \item \texttt{const const\_iterator} — константный объект типа \texttt{const\_iterator}.
\end{itemize}

\begin{remark}
\texttt{const iterator} и \texttt{const\_iterator} — это не одно и то же:
\begin{itemize}
    \item \texttt{iterator} / \texttt{const iterator} — речь о константности самого итераторного объекта;
    \item \texttt{const\_iterator} — речь о типе, который не позволяет менять элемент коллекции.
\end{itemize}
\end{remark}

\subsection{Почему для «левоконстантного» итератора требуется отдельный тип \texttt{const\_iterator}?}
Для указателей константность можно выразить прямо квалификатором \texttt{const}.  
Для итераторов этого недостаточно, потому что итератор — это уже не просто указатель, а абстрактный объект со своим интерфейсом.

Поэтому в стандартной библиотеке вводят отдельный тип:
\begin{itemize}
    \item \texttt{iterator} — для изменения элемента;
    \item \texttt{const\_iterator} — для только чтения элемента.
\end{itemize}

\begin{remark}
Это нужно, чтобы сам интерфейс итератора явно запрещал изменение элемента через «левоконстантный» доступ. Иначе был бы риск получить внешний объект, который выглядит как константный, но всё ещё позволяет менять данные через разыменование.
\end{remark}

\subsection{Краткие ответы на контрольные вопросы}
\begin{itemize}
    \item \textbf{Что такое косвенная адресация?} Доступ к объекту через промежуточную сущность.
    \item \textbf{Зачем она нужна?} Для передачи, обхода и абстрагирования доступа.
    \item \textbf{Какие механизмы используются?} Указатели, ссылки, итераторы, \texttt{string\_view}.
    \item \textbf{Сколько объектов обычно участвует?} Обычно два: косвенная сущность и объект-цель.
    \item \textbf{Что такое \texttt{const T*}?} Указатель на константный объект.
    \item \textbf{Что такое \texttt{T* const}?} Константный указатель.
    \item \textbf{Что такое \texttt{const T\&}?} Константная ссылка.
    \item \textbf{Что такое \texttt{const\_iterator}?} Итератор только для чтения элементов.
    \item \textbf{Зачем отдельный \texttt{const\_iterator}?} Чтобы запретить изменение элемента через итератор на уровне типа.
\end{itemize}
